% Gauss's law on an arbitrary closed surface S around a charge Q: E meets dA at
% an angle theta that varies from point to point, yet every field line crosses
% S, so the flux is still Q/epsilon_0. The companion figure,
% ch01-gauss-surface-sphere.tex, shows the same charge inside a sphere.
%
% Pulled in with \tikzfig{ch01-gauss-surface-arbitrary.tex}. A fragment, not a
% document -- one tikzpicture, no preamble -- so it inherits the fonts of
% whatever \inputs it. The charge style matches ch01-coulomb-law.tex, so Q
% looks the same in every figure of the chapter.
%
% Field lines and E are drawn so they cannot be mistaken for each other:
%   - a field line is a thin grey line from Q, with a small arrowhead partway
%     along to give its sense and none at its end -- it is a curve through
%     space, not a vector with a length;
%   - E is a black arrow that starts AT a point and whose length is its size.
%
% Every field line must cross S exactly once. That holds only while S is
% star-shaped as seen from Q -- keep it so if the blob is ever redrawn.
%
% The marked point on S is a joint of the to[out,in] path, where the tangent is
% the `out' angle by definition -- so the outward normal (to the right of the
% counter-clockwise travel) is exactly out - 90, with nothing to compute. E is
% aimed along (P) - (Q) with `let', so theta is whatever the geometry says it
% is, and the arc marking it follows if the blob is edited.
%
% Needs the arrows.meta and calc libraries, loaded in the preamble of
% gwbook.cls and _wrapper.tex -- see the note in ch01-coulomb-law.tex.
\begin{tikzpicture}[
    % The figsize knob: distances scale, lettering and node sizes do not.
    scale=1.3,
    >={Latex},
    charge/.style  = {circle, draw, fill=white, line width=0.9pt,
                      minimum size=12pt, inner sep=0pt, outer sep=2.5pt,
                      font=\normalsize},
    fieldline/.style = {line width=0.5pt, black!50},
    % The sense marker on a field line: small and narrow, so it reads as a
    % direction and not as the tip of a vector.
    sense/.style   = {-{Stealth[length=4pt, width=2.6pt]}},
    surface/.style = {dash pattern=on 4pt off 2pt, line width=1.1pt},
    % The curly pointer from a name to what it names: an open V of two strokes
    % (Straight Barb, 60 deg included angle) and a round cap, as in
    % cha-vt-diagram-main.tex -- a pen mark, not a vector.
    pointer/.style = {-{Straight Barb[round, length=4pt, width=4.6pt]},
                      line width=0.6pt, line cap=round, line join=round,
                      % Stop short of the target, so the tip points at it
                      % rather than touching it; shorten < does the same at
                      % the name end.
                      shorten >=3pt, shorten <=1pt},
    % Sans, so a note reads as an annotation rather than as part of the text
    % or the maths. (The S label sets the same font inline.)
    note/.style    = {font=\sffamily\scriptsize},
    % No white halo: no vector crosses a field line, and a halo would eat
    % into the base of whichever vector shares its starting point.
    vec/.style     = {->, line width=0.9pt},
    % The patch of surface that dA stands for.
    patch/.style   = {line width=2.6pt, line cap=butt},
    angle/.style   = {line width=0.6pt},
    % For labels that sit among the field lines.
    backed/.style  = {fill=white, inner sep=1.5pt},
  ]

  \pgfmathsetmacro{\lineR}{2.45}  % how far every field line runs from Q;
                                  % must clear S (max radius ~2.1)
  \pgfmathsetmacro{\nameR}{4.2}   % ... except the named 30 deg line, which
                                  % runs on to hold its name (about 2cm at the
                                  % book's 12pt) clear of S
  \pgfmathsetmacro{\senseR}{1.3}  % where along a line its sense marker sits:
                                  % inside S (min radius ~1.45, at the dent)
                                  % everywhere, and well clear of Q
  \pgfmathsetmacro{\alen}{1.15}   % length of dA
  \pgfmathsetmacro{\elen}{1.45}   % length of E
  \pgfmathsetmacro{\half}{0.16}   % half-width of the dA patch

  \node[charge] (Q) at (0,0) {$+$};
  % Two pieces per line: the first ends in the sense marker, the second
  % carries on from the marker's tip with no head of its own.
  \foreach \t in {0,30,...,330} {
    \pgfmathsetmacro{\len}{\t==30 ? \nameR : \lineR}
    \draw[fieldline, sense] (Q.\t) -- (\t:\senseR);
    \draw[fieldline] (\t:\senseR) -- (\t:\len);
  }

  % The field-line name is set along the 30 deg line itself, on the run past S
  % (which it crosses at radius ~1.95).
  \path (30:2.05) -- (30:\nameR)
    node[midway, sloped, above, inner sep=1.5pt, note] {electric field line};

  % Counter-clockwise, and smooth at every joint: each `out' is the previous
  % `in' + 180. The dent on top is what makes the shape read as arbitrary.
  \draw[surface]
    (1.93,-0.52) to[out=115, in=-20]
    (1.0,1.85)   to[out=160, in=0]
    (0.05,1.45)  to[out=180, in=-20]
    (-1.0,1.8)   to[out=160, in=90]
    % Named for the S pointer: a point between the 120 and 150 deg lines.
    node[pos=0.25, coordinate] (Sa) {}
    (-1.85,0.0)  to[out=270, in=160]
    (-0.8,-1.3)  to[out=-20, in=190]
    (0.6,-1.5)   to[out=10,  in=295]
    cycle;

  % Set above the tips of the 90 and 120 deg lines, which it would otherwise
  % run into.
  \node[anchor=base west] (Slbl) at (-2.9,2.75)
    {$S$ {\sffamily\scriptsize (Gaussian surface)}};
  \draw[pointer] ([xshift=4pt]Slbl.south west) to[out=-80, in=170] (Sa);

  \node[backed, anchor=north west] at (Q.-45) {$Q$};

  % The marked point is the (1.93,-0.52) joint, where the path starts:
  % tangent 115, normal 25. It lies at -15 deg from Q, midway between two
  % field lines, so E (which is radial) does not sit on top of one; and it is
  % on the right edge, so E, dA and theta open into the empty space there
  % rather than into the field lines.
  \coordinate (P) at (1.93,-0.52);
  \pgfmathsetmacro{\tang}{115}
  \pgfmathsetmacro{\nrm}{\tang-90}

  \draw[vec] (P) -- ++(\nrm:\alen) node[above right=-1pt] {$d\mathbf{A}$};
  \draw[vec] let \p1 = ($(P)-(Q)$), \n1 = {atan2(\y1,\x1)} in
    (P) -- ++(\n1:\elen) node[below right=-1pt] {$\mathbf{E}$};

  % theta, from dA round to E.
  \draw[angle] let \p1 = ($(P)-(Q)$), \n1 = {atan2(\y1,\x1)} in
    ($(P)+(\nrm:0.7)$) arc[start angle=\nrm, end angle=\n1, radius=0.7]
    ($(P)+({(\nrm+\n1)/2}:0.95)$) node {$\theta$};

  \draw[patch] ($(P)+(\tang:\half)$) -- ($(P)-(\tang:\half)$);

\end{tikzpicture}
